\chapter{Materials and Methods}

This chapter describes the methods developed and evaluated for automated artifact detection in fetal diffusion magnetic resonance imaging (dMRI). The research initially investigated three-dimensional brain segmentation and artifact classification approaches, followed by the development of a two-dimensional framework for brain localization and image-quality assessment. The following sections describe the imaging dataset, data preparation, segmentation and classification models, dynamic-level quality assessment, sequential processing, and experimental evaluation procedures.

\section{MRI Dataset and Ground-Truth Annotations}
\label{sec:mri-dataset-ground-truth}

The present research used low-field fetal diffusion magnetic resonance imaging (dMRI) data to develop and evaluate automated brain segmentation and artifact classification methods. The available dataset comprised volumetric diffusion-weighted MRI acquisitions, corresponding reference fetal-brain segmentation masks, and image-quality annotations assigned to diffusion dynamics. Initial experiments used complete three-dimensional (3D) volumes; subsequent experiments investigated two-dimensional (2D) slices. These approaches employed distinct prediction units and experimental datasets derived from the available imaging material.

\subsection{Dataset Source and MRI Acquisition}

The fetal dMRI data were provided through the research activities of the Smart Imaging Lab. The documented acquisition used a clinical 0.55\,T MAGNETOM Free.Max scanner (Siemens Healthineers, Forchheim, Germany), following the diffusion protocol described by Aviles Verdera et al.\ in the HERON framework \autocite{avilesverdera2025heron}.

The published protocol used single-shot diffusion-weighted spin-echo echo-planar imaging (DW-SE-EPI), with repetition time (TR) of 7200\,ms, echo time (TE) of 129\,ms, field of view (FOV) of $400\times400$\,mm$^2$, isotropic resolution of $3\times3\times3$\,mm$^3$, echo spacing of 0.97\,ms, and partial Fourier factor 6/8. Diffusion weighting included $b$-values of 0, 10, 50, 80, 200, 400, 600, and 1000\,s/mm$^2$, with three orthogonal diffusion directions for each non-zero $b$-value.

The published acquisition description reported 35 slices and an acquisition duration of approximately 2.43 minutes. The present project's image-processing records also expected a through-plane dimension of 36 slices for some volumes. These different recorded values are preserved rather than treated as equivalent without confirmation of the relevant imaging metadata. Because the models were developed on low-field acquisitions, their performance at different field strengths or with different acquisition protocols was not established.

\subsection{Dataset Organization and Diffusion Dynamics}

The imaging data were provided as Neuroimaging Informatics Technology Initiative (NIfTI) files and arranged in subject-specific directories containing individual diffusion acquisitions. In this thesis, a \emph{dynamic} denotes a three-dimensional diffusion-weighted image volume acquired under a particular diffusion-weighting configuration. Each dynamic comprises a spatial sequence of MRI slices.

The resulting data hierarchy comprises subjects, acquisition sequences, diffusion dynamics, and individual slices. The 3D approaches processed a whole dynamic to obtain one quality prediction. The 2D approach processed slices individually and investigated how their outputs could be combined into dynamic-level quality decisions.

Subject, dynamic, and slice identifiers were retained during data preparation to associate MRI inputs with reference brain masks and quality labels, maintain experimental partitions, and compare outputs generated by different stages. The supplied quality annotation table contained 38 case--sequence records with 22 dynamic annotation positions per record, resulting in 836 annotation entries. These are entries in the annotation inventory and should not automatically be equated to the number of available MRI image--mask pairs for any experiment.

\subsection{Reference Brain Segmentation Masks}

Reference fetal-brain segmentation masks accompanied the available volumetric MRI data. The masks were originally generated by an automatic procedure and subsequently manually corrected to refine anatomical delineation. For binary segmentation, foreground (1) represented the fetal brain and background (0) represented surrounding tissue and other image regions.

The reference masks served as voxel-wise targets for the volumetric brain segmentation experiments. For the 2D framework, matching image and mask slices were extracted along the same spatial axis and indexed by subject, dynamic, and slice position, preserving correspondence during supervised segmentation training and evaluation.

Reference masks also defined brain regions in experiments using reference-assisted, or \emph{oracle}, cropping. The word \emph{oracle} denotes use of the available reference mask rather than an automatically predicted mask; it does not establish perfect anatomical ground truth. Architecture-specific cropping procedures and predicted-mask comparisons are detailed later in this chapter.

\subsection{Artifact Quality Annotations}

Image-quality annotations were assigned to whole diffusion dynamics, not independently to individual slices. Three annotation categories were present: Good (0), Artifact (1), and Uncertain ($-1$). The original annotation inventory comprised 655 Good entries, 115 Artifact entries, and 66 Uncertain entries, for a total of 836. The corresponding proportions were approximately 78.3\%, 13.8\%, and 7.9\%.

\begin{table}[htbp]
\centering
\caption{Distribution of the original dynamic-level quality annotation entries. These counts refer to the annotation inventory, not to a model-specific evaluation cohort.}
\label{tab:original_dynamic_annotation_counts}
\begin{tabular}{lrr}
\hline
Quality annotation & Entries & Proportion \\
\hline
Good (0) & 655 & 78.3\% \\
Artifact (1) & 115 & 13.8\% \\
Uncertain ($-1$) & 66 & 7.9\% \\
\hline
Total & 836 & 100.0\% \\
\hline
\end{tabular}
\end{table}

Uncertain annotations were excluded when preparing binary artifact-classification targets, leaving 770 known-label annotation entries. This number describes label availability, rather than independently verifying the availability of 770 matched MRI--mask files.

The quality annotations describe complete dynamics, so they do not independently establish artifact visibility in each constituent slice. The 2D experiments used inherited labels, referred to as \emph{weak slice-level supervision}; their precise target assignment and implications for class imbalance are specified in Section~\ref{sec:data-preparation}.

\subsection{Dataset Composition and Selection Criteria}

Separate experimental populations were assembled for preliminary 3D segmentation, nnU-Net segmentation, 3D artifact classification, 2D brain segmentation, and 2D artifact classification. Availability of image--mask pairs, anatomical content, binary quality labels, and the intended prediction unit determined eligibility for each experiment. These populations were not interchangeable: the original annotation inventory does not describe the total number of usable images in any one analysis.

The 3D experiments used prepared volumetric samples, whereas the 2D pipeline used subject-separated slice populations. The segmentation dataset retained images with and without reference brain foreground, while the artifact classifier considered a restricted range of brain-containing slices. Exact entry counts, inclusion criteria, and split assignments are reported together in Section~\ref{sec:data-preparation}.

\section{Data Preparation and Experimental Partitioning}
\label{sec:data-preparation}

The volumetric and slice-wise model families required different image representations and preprocessing operations. Preparation included MRI--mask pairing, handling of invalid intensities, model-specific spatial normalization, partitioning of related observations, and construction of binary quality targets. The procedures below distinguish the preliminary 3D pipeline, the nnU-Net configuration, and the later 2D framework.

\subsection{MRI Volume Loading and Preprocessing}

Available NIfTI MRI volumes were associated with reference brain masks using case and dynamic identifiers. Spatial correspondence was checked at the image--mask pair level, and the reference masks were represented as binary labels distinguishing background from fetal-brain voxels. For 3D segmentation, each dynamic remained a complete volume with a single MRI input channel and a correspondingly shaped segmentation target.

The initial custom 3D U-Net used three-dimensional numerical tensors with an added channel dimension for volumetric convolution. In the subsequent nnU-Net experiment, each dynamic was assigned an individual case identifier and its paired image and reference mask were formatted according to the nnU-Net dataset specification, including the image-channel and segmentation-class definitions.

For 2D processing, the NIfTI volumes were read using NiBabel. The recorded extraction procedure identified a through-plane axis by an expected length of 36 when applicable; otherwise, the smallest dimension could be selected if it was unambiguous. MRI and mask arrays were sliced along the same identified axis, using corresponding slice indices. Each image slice was retained as a 32-bit floating-point array, and masks were binarized. Subject, dynamic, and slice identifiers were preserved for subsequent data selection, segmentation, classification, and dynamic-level aggregation.

\subsection{Intensity Normalization and Spatial Preparation}

Intensity preparation differed across the three experimental stages and was not assumed to follow one identical procedure.

\paragraph{Preliminary custom 3D U-Net.}
Non-finite MRI intensities were replaced with zeros, and each volume was cast to 32-bit floating-point representation. When the maximum voxel intensity was positive, the volume was divided by that maximum:
\begin{equation}
I_{\mathrm{norm}}(x,y,z)=\frac{I(x,y,z)}{I_{\max}},\qquad I_{\max}>0.
\end{equation}
The MRI volume and corresponding reference mask were then padded with zeros at the upper end of each spatial dimension to obtain dimensions divisible by eight. Identical spatial padding preserved the alignment of input and segmentation target. This preparation accommodated three successive downsampling operations without resizing the original volumetric anatomy.

\paragraph{nnU-Net v2.}
The nnU-Net experiments used the framework's dataset-specific automated planning and preprocessing rather than the custom U-Net's manual intensity scaling and padding. The resulting full-resolution model configuration, including target voxel spacing and patch size, is specified in Section~\ref{subsec:3d_nnunet_architecture}.

\paragraph{2D brain segmentation.}
Intensity normalization was performed independently for each extracted MRI slice. The 1st and 99th intensity percentiles were computed from finite, non-zero pixels when at least 32 such values were available; otherwise, all finite values were considered. Intensities were clipped to these bounds and scaled to $[0,1]$:
\begin{equation}
I_{\mathrm{norm}}=
\frac{\operatorname{clip}(I,P_1,P_{99})-P_1}{P_{99}-P_1},
\qquad P_{99}>P_1.
\end{equation}
The recorded implementation handled non-finite intensities and represented slices containing no finite intensities as zero arrays. For degenerate percentile bounds, the implementation-specific fallback must be verified rather than inferred from the equation.

The normalized MRI slices were resized to $160\times160$ pixels using bilinear interpolation. Reference masks were resized to matching dimensions with nearest-neighbor interpolation to preserve binary labels. Image and mask arrays were then represented as single-channel tensors for segmentation.

\paragraph{2D artifact classification.}
For classification, the pipeline generated brain-centered image crops from reference or predicted masks, using crop-specific normalization and resizing. The resulting $224\times224$ input representation and the different crop-generation variants are detailed in the developed 2D pipeline section.

\subsection{Subject-Level and Dynamic-Level Data Splits}

Distinct partitioning procedures were used during the earlier volumetric experiments and later 2D development. This matters because multiple dynamics belong to a subject and multiple correlated slices belong to a dynamic.

\paragraph{Preliminary custom 3D U-Net.}
A randomized 80:20 split of 770 image--mask entries, using seed 42, allocated 616 entries to training and 154 to testing. The split operated on volumetric entries, not independent subjects. Multiple dynamics from the same fetus could therefore appear in both partitions.

\paragraph{nnU-Net volumetric segmentation.}
The subsequent prepared dataset allocated 600 image--mask entries to training and 170 to testing, using a recorded random seed of 42. The corresponding manifest represented 35 subject identifiers, with subjects occurring in both partitions. Thus, these volume-entry evaluations must not be described as completely subject-independent. The later 3D classifier data contained 600 records, but their training, checkpoint-selection, and prediction-export relationship could not be established conclusively.

\paragraph{2D subject-level partitioning.}
The 2D framework used subject-based splitting to prevent slices and dynamics from the same fetus appearing across different partitions. A recorded two-stage stratified splitting procedure targeted approximately 70\% training, 15\% validation, and 15\% testing, using seed 42 and stratification by whether a subject included artifact-labelled observations. The documented realized allocation was 16 training subjects, four validation subjects, and four test subjects, corresponding to 24 subject groups. The available classification metadata confirmed disjoint subject identifiers across the three partitions, although the original split-generation procedure and complete mapping were not independently verified in all recovered records.

The prepared 2D segmentation inventories included both brain-present and background-only slices:
\begin{table}[htbp]
\centering
\caption{Prepared 2D segmentation inventories by subject-based partition and reference brain-mask presence.}
\label{tab:method_2d_segmentation_inventory}
\begin{tabular}{lrrr}
\hline
Partition & Total slices & Brain present & Brain absent \\
\hline
Training & 9,264 & 8,099 & 1,165 \\
Validation & 3,168 & 2,028 & 1,140 \\
Test & 3,168 & 1,902 & 1,266 \\
\hline
Total & 15,600 & 12,029 & 3,571 \\
\hline
\end{tabular}
\end{table}
A subset of background-only slices was included in segmentation training, whereas the validation and test sets retained background-only images for evaluation.

The 2D artifact-classification population was further limited using reference-mask brain presence and the slice-index range 8--25. The recorded implementation considered foreground area relative to the maximum area within each dynamic and excluded slices with insufficient brain content. The resulting dataset was:
\begin{table}[htbp]
\centering
\caption{Prepared 2D artifact-classification dataset by subject-based partition and inherited dynamic-level quality target.}
\label{tab:method_2d_classification_inventory}
\begin{tabular}{lrrr}
\hline
Partition & Good-labelled & Artifact-labelled & Total \\
\hline
Training & 4,996 & 1,082 & 6,078 \\
Validation & 1,405 & 101 & 1,506 \\
Test & 1,258 & 271 & 1,529 \\
\hline
Total & 7,659 & 1,454 & 9,113 \\
\hline
\end{tabular}
\end{table}
The selected slices retained their parent subject's partition assignment. The complete segmentation population and restricted classification population therefore represent different experimental analysis units and denominators.

\subsection{Handling Class Imbalance and Uncertain Labels}

Quality labels were available as Good (0), Artifact (1), and Uncertain ($-1$). Uncertain annotations were excluded from the construction of binary artifact-classification targets, rather than being assigned automatically to either class. Anatomical segmentation targets remained distinct from image-quality labels.

The binary quality annotations referred to full dynamics. For a slice $s$ extracted from dynamic $d$, the inherited target was
\begin{equation}
\widetilde{y}_{d,s}=Y_d,
\end{equation}
where $Y_d$ is the observed dynamic-level quality label. This weak supervision does not prove that every slice from an artifact-labelled dynamic visibly exhibits an artifact. Evaluation against the inherited targets therefore does not amount to evaluation against independent slice-level expert annotations.

The annotation inventory and the classification datasets were imbalanced towards Good observations. In the 2D classification partitions, artifact-labelled slices accounted for approximately 17.8\% of training data, 6.7\% of validation data, and 17.7\% of test data. This variation is important when interpreting precision, recall, and precision--recall measures.

An early EfficientNet-B0 oracle-crop classification baseline used class-weighted cross-entropy, with documented weights of 1.0 for Good and 4.6174 for Artifact. These weights approximately reflect the inverse relative prevalence of the training classes. They are specific to the recorded baseline: the available evidence does not establish that every subsequent oracle- and predicted-crop classifier used this identical weighting configuration. Model-specific optimization and checkpoint selection are described in Section 4.6.

Class imbalance and uncertainty in inherited targets are distinct issues. Weighting or sampling may modify the optimization contribution of minority-class examples but does not confirm that an inherited positive label corresponds to an artifact visible in an individual slice. For this reason, the subsequent evaluation reports class-sensitive metrics, identifies its observation unit, and distinguishes successful classifier predictions from missing inputs caused by segmentation or crop failure.

\section{Preliminary Three-Dimensional Model Development}
\label{sec:preliminary_3d}

The initial research phase investigated whether artifact quality could be assessed from complete three-dimensional fetal diffusion MRI dynamics by combining brain localization, anatomical cropping, handcrafted image descriptors, and conventional machine learning. This preliminary framework provided an exploratory alternative to learning image-quality features directly from the volume. The implementation and its limitations informed the subsequent development of an nnU-Net-based segmentation stage and a 3D residual neural network classifier.

\subsection{Initial 3D U-Net Brain Segmentation}
\label{subsec:initial_3d_unet}

The initial segmentation model was a custom three-dimensional U-Net that accepted single-channel volumetric MRI inputs and predicted binary fetal-brain masks. Its encoder comprised convolutional stages with 16, 32, and 64 feature channels, followed by a 128-channel bottleneck. Each stage used two $3\times3\times3$ convolutions with batch normalization and ReLU activations. Three $2\times2\times2$ max-pooling operations reduced spatial resolution. The decoder restored resolution using transposed convolutions and skip-connected encoder features, with decoder feature channels of 64, 32, and 16. A final $1\times1\times1$ convolution and sigmoid activation returned voxel-wise brain probabilities.

Training used the preliminary dynamic-entry partition described in Section~\ref{sec:data-preparation}. Because that partition was not subject-disjoint, its evaluation cannot be interpreted as an independent-subject test.

The model used the Adam optimizer with learning rate $10^{-4}$, batch size one, and 25 training epochs. Its reported objective added binary cross-entropy (BCE) to a Dice-based term:
\begin{equation}
\mathcal{L}=\mathcal{L}_{\mathrm{BCE}}+(1-\mathrm{DSC}_{\mathrm{thresholded}}).
\end{equation}
The implementation computed the Dice coefficient after thresholding the predicted probabilities at 0.5. Since this hard threshold did not provide a useful gradient with respect to the model outputs, the Dice term contributed to the numerical loss but not to gradient-based learning; parameter updates were effectively driven by BCE. The recorded training loop monitored average training loss and training Dice. A separate executed validation-based checkpoint-selection history was not established for this experiment.

During inference, sigmoid probabilities above 0.5 were converted to binary foreground masks. The masks were restored to the original image dimensions after removal of the spatial padding and passed to the brain-cropping procedure. Training curves and the limited segmentation outcomes are reported in Chapter 5 rather than used here as evidence of independent generalization.

\subsection{Volumetric Brain Cropping}
\label{subsec:preliminary_3d_cropping}

Predicted 3D segmentation masks were used to identify brain-centered regions of interest in the corresponding original MRI dynamics. For each nonempty mask, foreground voxel coordinates were collected and their minimum and maximum positions along the three spatial dimensions defined a rectangular bounding box. The cropping function accepted a configurable margin; the recorded inference call used an eight-voxel margin along each dimension, clipped to the image boundaries.

The cropped array retained original MRI intensities inside the bounding box rather than applying a voxel-wise binary masking operation. A crop could be generated only if the predicted mask contained foreground. The inference loop checked for empty masks and skipped crop saving when no brain region was detected.

Successful predicted masks and brain crops were saved in compressed NIfTI format. The predicted mask was saved with the original image affine, whereas the exploratory crop-saving code used an identity affine. Consequently, the saved cropped image was used as an intensity array for feature extraction and was not treated as preserving the source image's original world-coordinate mapping.

\subsection{Handcrafted Feature Extraction}
\label{subsec:preliminary_3d_features}

The exploratory classification approach represented each cropped MRI volume using five handcrafted statistical descriptors. These features were calculated directly from the positive-intensity voxels of the saved volumetric crop, using the \texttt{extract\_brain\_features} routine in the historical notebook.

Let $C$ denote a cropped MRI array and $V_+=\{C_i:C_i>0\}$ its positive-valued voxels. With $N_+=|V_+|$, the descriptors were:

\begin{enumerate}
\item \textbf{Mean positive intensity}: $\mu=N_+^{-1}\sum_{i=1}^{N_+}C_i$.
\item \textbf{Standard deviation}: $\sigma=\sqrt{N_+^{-1}\sum_{i=1}^{N_+}(C_i-\mu)^2}$, following NumPy's population-standard-deviation convention.
\item \textbf{Crop volume}: the number of voxels in the rectangular cropped array, $\prod_k D_k$, rather than physical brain volume.
\item \textbf{Nonzero voxel count}: the number of voxels satisfying $C_i>0$, recorded as \texttt{nonzero\_voxels}.
\item \textbf{Uniformity-labelled descriptor}: $U=\mu/(\sigma+10^{-5})$.
\end{enumerate}

The final descriptor was named \texttt{uniformity} in the implementation but is more precisely a stabilized mean-to-standard-deviation ratio; it should not be interpreted as a standardized image-texture uniformity measurement. If a crop contained no positive-valued voxels, the implementation returned zeros for all five descriptors.

The resulting five-dimensional representation,
\begin{equation}
\mathbf{x}=[\mu,\sigma,\prod_kD_k,N_+,U],
\end{equation}
was associated with the available dynamic-level quality target and passed to the preliminary Random Forest classifier. These simple descriptors summarized intensity and crop size but did not explicitly represent spatially localized artifacts or relationships between neighboring image regions.

\subsection{Random Forest Artifact Classification}
\label{subsec:preliminary_random_forest}

A Random Forest classifier was trained on the extracted five-feature vectors as an exploratory binary Good-versus-Artifact model. The scikit-learn implementation used 200 trees, maximum depth ten, balanced class weights, and random seed 42. Entries with uncertain ($-1$) quality labels were excluded.

The historical notebook converted the wide annotation table to long format and extracted the dynamic number from each crop filename. The crop-to-annotation join, however, matched only on the dynamic number and selected the first matching row. It did not also require the true subject or case identifier. Because a given dynamic number can occur in different subjects, this procedure did not establish a reliable correspondence between each cropped MRI input and its assigned ground-truth quality label.

The documented experiment contained only four resulting feature records. The splitting logic assigned a single observation to testing and the other three to training using seed 42. The Random Forest was then fitted to the training records and used to predict the one withheld label.

Because of the very small sample size and the incorrect crop-to-label matching rule, this experiment was retained as evidence that the feature-extraction and tabular-classification workflow was implemented, not as a valid held-out performance benchmark. Its output was not used to rank the Random Forest against the subsequent deep learning classifiers.

\subsection{Limitations of the Preliminary Approach}
\label{subsec:preliminary_3d_limitations}

The preliminary pipeline established an initial technical sequence from volumetric brain-mask prediction through region extraction and handcrafted feature classification. Nonetheless, it exposed limitations at several distinct stages.

First, the custom 3D U-Net inference inspection included volumes for which no foreground brain mask was detected, preventing downstream crop generation. The recorded optimization curves described training behavior but did not establish independently verified segmentation generalization. Second, the thresholded Dice term in the loss function did not provide the intended differentiable overlap-optimization signal. Third, the five global intensity and crop-size descriptors could not directly encode complex spatial artifact patterns. Finally, the Random Forest experiment combined very few records with an incorrect subject-to-label correspondence, making its reported predictive performance unsuitable for scientific comparison.

These observations motivated a more structured volumetric approach combining nnU-Net v2 fetal-brain segmentation with a 3D residual convolutional classifier that learned spatial representations directly from the brain-centered MRI volumes. The developed 3D pipeline is described in the following section.

\section{Developed 3D Artifact-Detection Pipeline}
\label{sec:developed_3d_pipeline}

Following the limitations identified in Section~\ref{subsec:preliminary_3d_limitations}, a developed volumetric pipeline combined nnU-Net v2 brain segmentation, anatomical cropping, and a three-dimensional residual convolutional neural network (3D ResNet) for artifact classification. Its processing design and offline acquisition-workflow simulation are described below.

\subsection{Overall 3D Pipeline Architecture}
\label{subsec:overall_3d_pipeline}

The framework connected volumetric fetal-brain localization and residual-network artifact classification through extraction of the predicted brain region. Each model's architecture and input preparation are specified in the dedicated subsections below.

The processing sequence can be written as
\begin{equation}
I_d \xrightarrow{f_{\mathrm{seg}}} M_d
\xrightarrow{g_{\mathrm{ROI}}} C_d
\xrightarrow{f_{\mathrm{cls}}} \hat{Y}_d,
\end{equation}
where $I_d$ is the input MRI dynamic, $f_{\mathrm{seg}}$ is the trained segmentation model, $M_d$ is its predicted brain mask, $g_{\mathrm{ROI}}$ represents brain-region extraction and preprocessing, $C_d$ is the resulting crop, $f_{\mathrm{cls}}$ is the volumetric classifier, and $\hat{Y}_d$ is the predicted quality class.

The segmentation stage was intended to restrict classification to the anatomical region of interest rather than the entire fetal and maternal imaging field. Because the classifier depended on the predicted mask, however, errors in brain localization could propagate into cropping and classification. Each successfully prepared dynamic produced one quality prediction; aggregation of independent slice-level outputs was not required in this volumetric architecture.

\subsection{nnU-Net v2 Fetal Brain Segmentation}
\label{subsec:3d_nnunet_architecture}

nnU-Net v2 was used for volumetric fetal-brain localization, with configuration determined through dataset-specific framework planning.

\paragraph{Dataset preparation.}
Each diffusion dynamic was represented as a single-channel three-dimensional MRI input with a paired binary brain mask. Files were associated through unique case identifiers in the nnU-Net dataset specification. The prepared segmentation population, its training/test allocations, and the volume-entry partitioning limitation are documented in Section~\ref{sec:data-preparation}.

\paragraph{Network configuration.}
The reported framework configuration was \texttt{3d\_fullres}, with target voxel spacing $3\times3\times3$\,mm$^3$ and a patch size of $40\times160\times160$ voxels. The documented six-stage encoder--decoder used feature channels of 32, 64, 128, 256, 320, and 320. Three-dimensional convolutions were combined with instance normalization and leaky rectified linear unit (LeakyReLU) activations. The encoder learned multiscale volumetric features, while the decoder restored spatial resolution, aided by skip connections carrying corresponding encoder information.

\paragraph{Segmentation inference.}
Each dynamic was processed independently with the trained nnU-Net v2 model. The documented inference used the model associated with fold 0. The segmentation stage generated the predicted brain mask,
\begin{equation}
M_d=f_{\mathrm{nnU\text{-}Net}}(I_d),
\end{equation}
which supplied the anatomical foreground coordinates for downstream cropping. The segmentation model performed localization, not direct image-quality classification.

\subsection{Three-Dimensional Brain Region Extraction}
\label{subsec:3d_roi_extraction}

The predicted mask determined the bounding box used to extract the original brain-centered MRI intensities for classification.

\paragraph{Bounding-box construction.}
The minimum and maximum coordinates of the mask's foreground voxels defined a bounding box in each spatial dimension. The documented cropping procedure expanded this box by two voxels along each axis, clipped to the original image boundaries. Conceptually, for dimension $k$ with image size $D_k$ and foreground-coordinate set $\Omega_k$, the lower and upper crop positions were defined by
\begin{equation}
a_k=\max(0,\min(\Omega_k)-m),\qquad
b_k=\min(D_k,\max(\Omega_k)+m),\qquad m=2.
\end{equation}
The exact handling of the upper array-slicing endpoint depends on the implementation. Bounding-box cropping retained intensity information inside the rectangular crop rather than multiplying the original image by the binary foreground mask.

\paragraph{Intensity normalization.}
The extracted region was normalized using its 1st and 99th intensity percentiles. Values outside the percentile interval were clipped, and the remaining intensities were mapped to $[0,1]$. For valid percentile bounds, this is represented by
\begin{equation}
C_{\mathrm{norm}}=
\frac{\operatorname{clip}(C,P_1,P_{99})-P_1}
{P_{99}-P_1},\qquad P_{99}>P_1,
\end{equation}
where $C$ is the cropped volume. The normalized crop was converted into a single-channel volumetric classifier input.

\paragraph{Segmentation-failure handling.}
A nonempty foreground mask was required to define a bounding box. If the predicted segmentation contained no foreground, the crop could not be generated in the usual manner. Archived implementations included explicit handling for such cases, although their responses were not identical across versions. Furthermore, a nonempty mask permitted crop generation without guaranteeing accurate anatomical localization. Crop availability and anatomical correctness are therefore distinct properties.

\subsection{3D ResNet Artifact Classification Architecture}
\label{subsec:3d_resnet_architecture}

The downstream classifier was a custom three-dimensional residual convolutional neural network designed to distinguish Good and Artifact dynamics from brain-centered volumetric MRI inputs. In contrast to the preliminary five-feature Random Forest approach, the ResNet learned spatial image representations through convolutional feature extraction.

\paragraph{Initial convolution.}
A single-channel input crop passed through a $3\times3\times3$ convolution producing 32 feature channels, followed by batch normalization and ReLU activation.

\paragraph{Residual stages.}
Three residual stages increased feature channels to 64, 128, and 256 while reducing spatial resolution through strided convolutions. Each residual block contained two $3\times3\times3$ convolutions with batch normalization, alongside a shortcut connection. The general residual operation can be written as
\begin{equation}
\mathbf{y}=F(\mathbf{x},\theta)+S(\mathbf{x}),
\end{equation}
where $F$ denotes the convolutional transformation and $S$ the shortcut pathway. When dimensions differed, a $1\times1\times1$ convolution adjusted the shortcut to the output shape.

\paragraph{Pooling and classification head.}
Adaptive average pooling reduced the output feature maps to a 256-dimensional representation. A fully connected layer projected this to 128 hidden units, followed by ReLU activation and dropout with probability 0.5. The output layer produced two logits corresponding to Good and Artifact.

For an input crop $C_d$, let the output logits be $\mathbf{z}_d=f_{\mathrm{ResNet}}(C_d)$. Softmax converted them to class probabilities,
\begin{equation}
p_{d,k}=\frac{\exp(z_{d,k})}{\sum_{j=0}^{1}\exp(z_{d,j})},
\end{equation}
where $p_{d,1}$ denotes the artifact-class score for dynamic $d$.

\subsection{Training Configuration and Model Selection}
\label{subsec:3d_training_model_selection}

The segmentation and classification stages were trained for different objectives: anatomical fetal-brain localization and binary image-quality prediction, respectively.

\paragraph{nnU-Net segmentation training.}
The nnU-Net v2 model used the paired volumes and reference masks described in Section~\ref{subsec:3d_nnunet_architecture}. The recovered materials did not include a complete execution-linked training history or independently verified aggregate held-out segmentation results. Consequently, its precise selected checkpoint epoch and generalization performance were not inferred from the preparation records.

\paragraph{3D ResNet classification data.}
The recovered classifier dataset comprised 600 eligible dynamic-level records from 35 subjects, with 508 Good and 92 Artifact labels (15.3\% Artifact). Unlike the larger segmentation dataset, these entries were associated with prepared brain-centered classifier inputs. The relationship between the executed training split and prediction export was not completely recoverable.

\paragraph{Classification optimization.}
The available 3D ResNet training implementation specified 40 epochs, weighted sampling for class imbalance, and a focal-loss objective. Weighted sampling adjusted representation of the two classes during training, while focal loss placed relatively greater emphasis on difficult classification examples. The general loss family can be expressed as
\begin{equation}
\mathcal{L}_{\mathrm{focal}}=-\alpha_t(1-p_t)^\gamma\log(p_t),
\end{equation}
where $p_t$ is the predicted probability of the correct class, $\alpha_t$ is a weighting parameter, and $\gamma$ controls emphasis on difficult examples. This formula does not independently establish the exact hyperparameters used to produce each archived classifier checkpoint.

\paragraph{Partitioning and model selection.}
The documented classifier training implementation used stratified dynamic-record partitioning rather than subject-grouped splitting, as distinguished from the 2D partitions in Section~\ref{sec:data-preparation}. The training script specified checkpoint selection by the artifact-class F1-score. The recovered records do not conclusively link the executed training history, selected epoch, checkpoint, and later prediction export; the implications for evaluation independence are addressed with the classification results.

\subsection{Dynamic-Level Artifact Prediction}
\label{subsec:3d_dynamic_prediction}

The developed 3D classifier predicted image quality directly for each complete diffusion dynamic. After segmentation, cropping, and normalization, a usable brain-centered volume was passed to the 3D ResNet, yielding two class scores and corresponding softmax probabilities. The predicted class was selected as
\begin{equation}
\hat{Y}_d=\arg\max_{k\in\{0,1\}}p_{d,k},
\end{equation}
where class 0 denotes Good and class 1 denotes Artifact. The artifact-class probability was retained as a continuous score.

Unlike the later 2D approach, this dynamic-level classification did not require aggregation of independently predicted slice scores. Nonetheless, a prediction depended on the availability of the source MRI volume, a usable predicted segmentation mask, and a successfully generated brain crop.

Predictions were available for the eligible downstream classification population described in Section~\ref{subsec:3d_training_model_selection}. Since no complete execution log traced all candidate acquisitions through segmentation, cropping, and prediction, the export alone could not establish an end-to-end pipeline completion rate.

\subsection{Offline Acquisition-Workflow Simulation}
\label{subsec:3d_offline_workflow}

An additional development experiment investigated the integration of the volumetric stages in an offline, acquisition-oriented quality-assessment workflow. The aim was to examine how newly available MRI files could automatically trigger model processing instead of requiring separate manual execution of segmentation, cropping, and classification.

\paragraph{File-triggered workflow.}
The historical implementation described a folder-monitoring mechanism that detected newly saved NIfTI volumes and initiated MRI loading, fetal-brain segmentation, volumetric ROI extraction, artifact classification, and output logging. The simulation approximated the processing that would follow availability of a completed diffusion dynamic.

\paragraph{Logging and demonstration interface.}
Historical reports described recording file arrival, processing initiation, processing completion, and the predicted quality category. A Streamlit-based upload-and-predict interface was also presented as an interactive demonstration of volumetric quality assessment.

\paragraph{Computational motivation.}
Historical project updates identified complete-volume segmentation as a substantial contribution to overall processing time. The timings were obtained at different stages of development without fully documented identical hardware and measurement boundaries. Numerical timing results and comparisons with the 2D system are consequently presented separately in Chapter 5.

\paragraph{Scope and limitation.}
The simulation processed MRI files after they became available; it did not establish inference during active MRI acquisition or communication with the scanner acquisition controller. Neither automatic interruption of an ongoing scan nor reacquisition of a corrupted dynamic was demonstrated. The experiment therefore provides evidence of an offline file-triggered processing workflow, not prospective scanner-integrated real-time operation. Its dependence on complete-volume availability and the associated computational burden motivated the investigation of slice-wise quality assessment described in Section 4.5.

\section{Motivation for Slice-Wise Artifact Detection}
\label{sec:motivation_2d}

The developed volumetric framework combined nnU-Net v2 fetal-brain localization and 3D ResNet artifact classification to produce one quality prediction per complete diffusion MRI dynamic. Preserving three-dimensional context was useful for representing spatial image features, but the wider research objective also required consideration of when quality feedback could become available. This motivated investigation of slice-wise processing as a distinct computational and decision-making strategy, not an assumption that two-dimensional images necessarily provided superior artifact information.

\subsection{Limitations of Complete-Volume Processing for Early Feedback}
\label{subsec:3d_feedback_limitations}

\paragraph{Dependence on complete-volume availability.}
The 3D segmentation model accepted a reconstructed volumetric input. The resulting brain mask determined a three-dimensional region of interest, which was then used to prepare the classifier input. Consequently, the ordinary 3D processing sequence required a complete usable dynamic before segmentation, cropping, and classification could yield a quality decision:
\begin{equation}
I_d\longrightarrow M_d\longrightarrow C_d\longrightarrow p_d,
\end{equation}
where $I_d$ is the input dynamic, $M_d$ its predicted brain mask, $C_d$ the brain-centered volume, and $p_d$ the classifier's artifact probability. The design permitted retrospective dynamic-level assessment but did not provide a mechanism to evaluate an individual acquired image section using that same volumetric model.

\paragraph{Processing cost and dependency between stages.}
The developed framework required volumetric segmentation before the 3D ResNet could classify a brain crop. The historical offline scanner-workflow simulation identified nnU-Net processing as an important bottleneck. Development reports described approximately 15--25 seconds per scan in June 2026 and approximately 17--20 seconds per volume in a later presentation. These were historical observations obtained under incompletely matched conditions; their hardware, checkpoints, input preparation, and timing boundaries could not be independently established as identical. Accordingly, the values motivated further investigation without being interpreted as a definitive runtime benchmark for the final 3D pipeline. A separate prepared-crop 3D ResNet benchmark excluded nnU-Net segmentation and initial crop preparation, so it could not substitute for full-pipeline timing.

The dependence on brain localization also introduced a potential failure point. An empty segmentation mask prevented the usual bounding-box crop, while an anatomically inaccurate but nonempty mask could still generate an unsuitable input. These limitations were relevant to both the original 3D pipeline and the subsequent 2D design; changing dimensionality alone did not remove them.

\paragraph{Implications for acquisition-oriented quality control.}
The file-watching demonstration described in Section~\ref{subsec:3d_offline_workflow} triggered processing when newly saved NIfTI volumes became available. It did not demonstrate direct scanner communication, analysis during active acquisition, or automatic interruption and reacquisition of corrupted dynamics. Thus, the feasibility of retrospectively processing a completed acquisition remained distinct from the ability to supply feedback before a dynamic was finished. The desire to investigate smaller input units, lower computational overhead, and potentially earlier decisions motivated the slice-wise approach.

\subsection{Transition from 3D Dynamics to 2D Slices}
\label{subsec:transition_3d_to_2d}

The next development stage treated the individual two-dimensional MRI slice as the segmentation and classification input. For dynamic $d$ containing $S_d$ ordered slices, let
\begin{equation}
I_d=\{I_{d,1},I_{d,2},\ldots,I_{d,S_d}\}.
\end{equation}
The 2D framework generated a classifier artifact score $p_{d,s}$ for each successfully processed eligible slice $I_{d,s}$, preserving the slice's subject and dynamic identity. Unlike the 3D classifier, it did not require a complete volumetric tensor as the input to either of its neural networks. The ordered slice scores could then be examined individually or combined into a dynamic-level quality decision.

\paragraph{Preserving anatomical localization.}
The transition changed the image representation without abandoning the segmentation-to-classification architecture. A MobileNetV3-Large-based U-Net generated brain masks from individual MRI slices. These masks provided bounding boxes for brain-centered image crops, which were subsequently classified using EfficientNet-B0. The trained components, cropping methods, and inference procedures are detailed in Section~\ref{sec:developed_2d_pipeline}.

\paragraph{Potential computational and decision advantages.}
Processing smaller 2D inputs enabled investigation of efficient neural network architectures and independent slice inference. Retention of the original slice order also made it possible to formulate a sequential stopping rule: an alert could potentially be generated after a subset of slices if the available scores met a selected condition. This was a design opportunity, not evidence that the implemented system had already achieved prospective early artifact detection. The recorded experiments operated on slices extracted from available NIfTI data.

\paragraph{Trade-offs of losing volumetric context.}
A 2D classifier did not directly observe spatial relationships across adjacent slices. Artifacts requiring full volumetric context could therefore be more difficult to identify from individual images. Further, source quality labels were assigned to whole dynamics. For eligible slice $s$ belonging to dynamic $d$, the classification target inherited the original label, $\widetilde y_{d,s}=Y_d$. This weak supervision did not establish that the visible artifact was present in each positive-labelled slice. Finally, segmentation and crop failures remained possible at the slice level. These issues required explicit evaluation of slice classification, crop reliability, and dynamic-level decision rules, rather than equating rapid single-slice inference with reliable acquisition-quality assessment.

\subsection{Design Requirements for the 2D Framework}
\label{subsec:2d_design_requirements}

The transition motivated a set of requirements guiding the developed 2D framework.

\paragraph{Reliable fetal-brain localization.}
Each image slice required a predicted binary brain mask so that classification could focus on the relevant anatomical region. The segmentation approach needed to account for slices with small, absent, or difficult-to-localize brain foreground. A MobileNetV3-Large U-Net was selected as the developed segmentation architecture.

\paragraph{Computationally efficient processing.}
The segmentation and classification stages were designed for individual two-dimensional inputs, allowing their inference costs to be benchmarked separately. The runtime of preprocessing, segmentation, crop preparation, classification, and decision logic was measured rather than inferred solely from architecture size. Stage-specific measurements did not, by themselves, constitute total scanner-to-decision latency.

\paragraph{Anatomically focused classification and crop evaluation.}
The EfficientNet-B0 classifier received intensity crops centered on the predicted fetal brain. Reference-mask-assisted oracle crops were also prepared to study alternative localization conditions. Because separately trained oracle and predicted-crop classifiers were used in documented experiments, any differences in their results could not be attributed exclusively to segmentation quality.

\paragraph{Identity preservation and acquisition-level decisions.}
Subject, dynamic, and slice identifiers were retained throughout processing. They supported subject-separated model evaluation, inheritance of dynamic-level quality targets, grouping of slice probabilities, and comparison with acquisition-level annotations. The overall quality-control framework required an explicit rule for translating multiple slice results into one dynamic-level decision.

\paragraph{Explicit processing-failure handling.}
An unusable or empty predicted brain mask was treated as an unavailable classifier input, not automatically as a low artifact probability. Segmentation failures were recorded separately and could be incorporated into dynamic-level quality-control rules. A failure signal was not assumed to be independently confirmed evidence of visible artifact corruption.

\paragraph{Sequential assessment with clearly bounded claims.}
The architecture needed to permit evaluation of slices in order and optional stopping after a classifier or localization-failure trigger. A demonstrated stopping rule must be distinguished from measured first-detection latency and prospective scanner integration. The implemented aggregation and sequential-decision procedures are specified in Section~\ref{sec:dynamic_sequential_assessment}.

Taken together, these requirements motivated a lightweight slice-wise segmentation-to-classification workflow and a separate dynamic-level decision layer. The transition was therefore driven by investigation of feedback timing and processing strategy while retaining the known limitations of weak labels, segmentation dependence, and incomplete volumetric context.

\section{Developed 2D Artifact-Detection Pipeline}
\label{sec:developed_2d_pipeline}

Following the developed volumetric approach, the two-dimensional framework used MobileNetV3-Large U-Net for fetal-brain localization and EfficientNet-B0 for artifact classification of brain-centered intensity crops. Individual slices were processed independently, with subject, dynamic, and slice identifiers retained to support subsequent aggregation. General dataset preparation and splits are established in Section~\ref{sec:data-preparation}.

\subsection{Overall 2D Pipeline Architecture}
\label{subsec:2d_overall_architecture}

The framework connected two neural network stages through anatomical cropping. For input MRI slice $I_{d,s}$ belonging to dynamic $d$, processing was represented by
\begin{equation}
I_{d,s}\xrightarrow{f_{\mathrm{seg}}}\hat M_{d,s}
\xrightarrow{g_{\mathrm{ROI}}}C_{d,s}
\xrightarrow{f_{\mathrm{cls}}}p_{d,s},
\end{equation}
where $\hat M_{d,s}$ is the predicted brain mask, $C_{d,s}$ the cropped original image intensities, and $p_{d,s}$ the artifact-class score. Segmentation masks determined crop location and were not additional classifier image channels. Each model accepted 2D inputs, with independent slice predictions retained for dynamic-level decisions.

\subsection{Slice Extraction and Preprocessing}
\label{subsec:2d_slice_preprocessing}

\label{subsubsec:2d-normalization}
\label{subsubsec:artifact-labeling}

NIfTI MRI volumes and corresponding masks were sliced along the same axis using NiBabel, preserving subject, dynamic, and slice identifiers. The recorded axis-selection procedure used an expected dimension of 36 where applicable and otherwise a uniquely identifiable smallest spatial dimension. Extracted images were represented as 32-bit floating-point arrays; masks were binarized.

For segmentation, each slice underwent percentile-based normalization and resizing to $160\times160$ pixels, with bilinear interpolation for MRI intensities and nearest-neighbor interpolation for reference masks. For classification, eligible brain-containing slices were identified from reference-mask foreground area relative to the maximum within each dynamic and restricted to slice indices 8--25. Their Good (0) or Artifact (1) targets were inherited from parent-dynamic annotations, while uncertain ($-1$) labels were excluded. Detailed normalization, data counts, and partitioning are described in Section~\ref{sec:data-preparation}. The inherited targets are weak slice-level supervision, not independent artifact visibility labels.

\subsection{MobileNetV3-U-Net Segmentation Architecture}
\label{subsec:2d_mobilenet_architecture}

A two-dimensional segmentation network combining MobileNetV3-Large with a U-Net-based decoder was developed for automated fetal brain localization in diffusion MRI slices. The architecture employed MobileNetV3-Large as the encoder to extract hierarchical image features, while a custom decoder reconstructed the spatial information required for pixel-wise brain segmentation. The network processed single-channel grayscale MRI images of $160 \times 160$ pixels and produced a corresponding segmentation map identifying the fetal brain region. The proposed architecture is illustrated in Figure~\ref{fig:mobilenetv3_unet_architecture}.

\begin{figure}[htbp]
\centering
\includesvg[width=\textwidth]{mobilenetv3_unet_architecture}
\caption{Architecture of the MobileNetV3-Large U-Net for two-dimensional fetal brain segmentation. Five encoder feature scales connect to the four-stage decoder through skip connections. A final convolution produces segmentation logits, which are converted to a binary mask using sigmoid activation and a threshold of 0.5.}
\label{fig:mobilenetv3_unet_architecture}
\end{figure}

The encoder was initialized using the ImageNet-pretrained MobileNetV3-Large architecture, excluding its original classification head. Since the pretrained network was designed for three-channel RGB images, the first convolutional layer was modified to accept a single grayscale input channel. The pretrained convolutional weights were adapted by averaging the weights across the three original input channels. Feature maps were extracted at five spatial resolutions, with the deepest representation serving as the bottleneck and the remaining four representations providing skip connections to the decoder.

The decoder consisted of four successive upsampling blocks with output channels of 256, 128, 64, and 32. At each stage, the decoder feature maps were upsampled using bilinear interpolation to match the spatial dimensions of the corresponding encoder features. The upsampled features were concatenated with the encoder features through skip connections and processed using two consecutive $3 \times 3$ convolutions, each followed by batch normalization and the Hardswish activation function. Following the final decoder block, the feature maps were upsampled to the original input resolution and passed through an additional convolutional block containing 24 channels.

A final $1 \times 1$ convolution reduced the feature representation to a single-channel output containing pixel-wise segmentation logits. During inference, a sigmoid activation converted these logits into probabilities between 0 and 1, representing the likelihood of each pixel belonging to the fetal brain. The resulting probabilities were thresholded at 0.5 to obtain binary segmentation masks. The developed MobileNetV3-Large U-Net architecture contained 6,690,633 trainable parameters.

The complete single-slice segmentation procedure is illustrated in Figure~\ref{fig:single_slice_segmentation_workflow}. Each extracted MRI slice was normalized and resized to $160 \times 160$ pixels before being passed through the segmentation network. The resulting binary mask provided the spatial information required for subsequent brain-centered region-of-interest extraction.

\begin{figure}[htbp]
\centering
\includesvg[width=\textwidth]{single_slice_segmentation_workflow}
\caption{Single-slice segmentation workflow. Two-dimensional fetal MRI slices undergo intensity normalization and resizing before inference with the MobileNetV3-Large U-Net. The predicted brain masks support subsequent region-of-interest localization.}
\label{fig:single_slice_segmentation_workflow}
\end{figure}

\subsection{Segmentation Training and Model Selection}
\label{subsec:2d_segmentation_training}

The MobileNetV3-Large U-Net was trained using two-dimensional fetal diffusion MRI slices and their corresponding reference brain masks. The training dataset included slices containing fetal brain tissue and a subset of background-only slices to improve the network's ability to distinguish the brain from surrounding maternal and fetal anatomy. All brain-containing training slices were retained, while 25\% of the available background-only slices were sampled for inclusion. The subject-independent validation dataset was used to monitor segmentation performance and select the best-performing model.

The segmentation network was optimized using a combined Dice and binary cross-entropy (BCE) loss function. Dice loss measured the spatial overlap between predicted brain regions and reference segmentation masks, whereas BCE loss penalized errors in individual pixel predictions. The combined objective was defined as
\begin{equation}
\mathcal{L}_{\mathrm{seg}} =
0.7\mathcal{L}_{\mathrm{Dice}}+
0.3\mathcal{L}_{\mathrm{BCE}}.
\end{equation}
The Dice component was calculated using sigmoid-transformed predictions with a smoothing constant of 1 to maintain numerical stability. Binary cross-entropy was computed directly from the segmentation logits. A higher weight was assigned to Dice loss to emphasize accurate foreground overlap, particularly because the fetal brain occupied only a portion of the complete MRI image.

Data augmentation was applied exclusively to the training images to improve robustness to variations in fetal orientation, position, and MRI signal intensity. Spatial augmentation consisted of horizontal flipping with a probability of 0.5, random rotations between $-8^\circ$ and $+8^\circ$, translations between $-0.04$ and $+0.04$ in normalized affine-grid coordinates, and scaling factors ranging from 0.95 to 1.05. Identical geometric transformations were applied to each MRI slice and its corresponding segmentation mask. Bilinear interpolation was used for image transformations, while nearest-neighbor interpolation preserved the binary mask labels. Intensity augmentation included random gain adjustments between 0.90 and 1.10 and intensity offsets between $-0.05$ and $+0.05$, applied with a probability of 0.7. Gaussian noise with a standard deviation ranging from 0 to 0.025 was additionally introduced with a probability of 0.3. Modified image intensities were constrained to the interval $[0,1]$. No random augmentation was applied to the validation or test datasets.

The network was trained using the AdamW optimizer with an initial learning rate of $1 \times 10^{-3}$, weight decay of $1 \times 10^{-4}$, and a batch size of 32. The training configuration specified a maximum of 80 epochs and a random seed of 42. A ReduceLROnPlateau scheduler monitored validation Dice and reduced the learning rate by a factor of 0.5 after four epochs without improvement, subject to a minimum learning rate of $1 \times 10^{-6}$. Gradient clipping with a maximum norm of 5 was applied to improve training stability. Automatic mixed-precision computation was enabled when training on a compatible CUDA GPU.

After each epoch, segmentation performance was assessed on the validation dataset using the Dice similarity coefficient. The model parameters were saved whenever validation Dice improved, and early stopping was applied after 12 consecutive epochs without improvement. The checkpoint achieving the highest validation Dice was selected for subsequent inference and evaluation on the held-out test dataset.

\paragraph{Segmentation inference.}
Following network training, the checkpoint achieving the highest validation Dice similarity coefficient was selected for fetal brain segmentation on unseen two-dimensional diffusion MRI slices. During inference, the trained MobileNetV3-Large U-Net was initialized with its saved model parameters and placed in evaluation mode. Predictions were generated without gradient computation to reduce computational overhead.

Each MRI slice underwent the same intensity normalization and spatial resizing procedures described in Section~\ref{subsec:2d_slice_preprocessing}. The normalized image was resized to $160 \times 160$ pixels and represented as a single-channel tensor before being passed through the segmentation network. The network produced a pixel-wise logit map, which was converted into foreground probabilities using the sigmoid activation function:
\begin{equation}
P(x,y)=\frac{1}{1+e^{-z(x,y)}},
\end{equation}
where $z(x,y)$ denotes the predicted logit at pixel location $(x,y)$, and $P(x,y)$ represents the corresponding probability of belonging to the fetal brain.

A fixed probability threshold of 0.5 was applied to convert the predicted probability map into a binary segmentation mask:
\begin{equation}
\hat{M}(x,y)=
\begin{cases}
1, & P(x,y)\geq 0.5,\\
0, & P(x,y)<0.5.
\end{cases}
\end{equation}
Pixels exceeding the threshold were classified as fetal brain foreground, while the remaining pixels were assigned to the background. For subsequent processing of the original MRI slices, the predicted binary masks were resized from $160 \times 160$ pixels to the original image dimensions using nearest-neighbor interpolation. Thresholding was performed before this spatial restoration to preserve the binary representation of the predicted masks.

The resulting segmentation masks were used to localize the fetal brain independently within each MRI slice. Foreground pixels provided the spatial coordinates required for brain bounding-box extraction and subsequent region-of-interest cropping. When a predicted mask contained no foreground pixels, the brain region could not be localized and the corresponding slice was identified as a segmentation failure. These failures were recorded separately from successfully segmented slices and considered during downstream artifact classification and dynamic-level quality assessment.

\subsection{Brain Bounding-Box Extraction and Cropping}
\label{subsec:2d_roi_crop}

\label{subsubsec:roi-cropping}

The binary fetal brain masks obtained from two-dimensional segmentation were used to extract brain-centered regions of interest (ROIs) from the corresponding original diffusion MRI slices. Cropping retained the image intensities in the selected rectangular region, including surrounding tissue within the crop, rather than multiplying the image by the segmentation mask. This reduced the field of view presented to the artifact classifier while preserving local anatomical context.

For each nonempty mask, the minimum and maximum foreground row and column coordinates defined the tight brain bounding box. The upper coordinates were made exclusive for array extraction. If its height and width were $h$ and $w$, respectively, the margins on each side were defined by
\begin{equation}
m_y=\left\lceil0.15h\right\rceil,\qquad
m_x=\left\lceil0.15w\right\rceil.
\end{equation}
The bounding box was expanded in both dimensions by these margins. A square crop was then constructed around the expanded box center, with side length equal to the larger expanded dimension. Coordinates falling outside the original slice were handled by extracting the available pixels and zero-padding the remainder of the square.

The resulting intensity crop underwent the normalization operation implemented in the crop-generation procedure, followed by bilinear resizing to $224\times224$ pixels (align\_corners=False). The resulting single-channel floating-point representation was provided to the artifact classifier. A mask without foreground pixels did not yield a valid bounding box and was recorded as an unsuccessful localization rather than replaced by an arbitrary crop.

\subsection{Oracle and Predicted Mask-Based Crops}
\label{subsec:2d_oracle_predicted_crops}

Two ROI-generation conditions were used to investigate the effect of brain localization on subsequent classification. In the oracle condition, the reference fetal brain segmentation mask supplied the bounding-box coordinates. In the predicted condition, the mask generated by the trained MobileNetV3-Large U-Net supplied the coordinates. Both conditions used the same square-crop construction, margin, padding, normalization, and resizing procedure.

The term \emph{oracle} denotes reference-mask-assisted localization, not perfectly accurate ground truth. The reference masks originated from automated segmentation followed by manual correction. For the predicted condition, errors in the estimated mask could change the crop center, extent, and surrounding image content even when some brain foreground was detected. Empty predicted masks were treated as crop-generation failures; a nonempty but inaccurate mask could still produce a crop.

The original slice identity, dynamic identifier, and subject identity were retained so that corresponding oracle and predicted crops could be compared. In the documented test classification population, 1,529 slices were eligible for reference-mask crops, whereas 1,499 yielded predicted-mask crops and 30 had no usable predicted crop. These counts describe input availability; they do not establish the classification accuracy of either cropping condition.

\subsection{EfficientNet-B0 Classification Architecture}
\label{subsec:2d_efficientnet_architecture}

Artifact classification was performed using EfficientNet-B0, a compact convolutional neural network adapted for two-class prediction from brain-centered fetal diffusion MRI slices. The model accepted a single grayscale crop of $224\times224$ pixels. The input images were represented as floating-point arrays, with nonfinite values replaced and intensity values constrained to $[0,1]$ before model inference.

The implementation used the EfficientNet-B0 feature extractor from torchvision. Where pretrained initialization was requested, the ImageNet-1K pretrained weights were loaded. The first convolution was replaced with a one-channel convolution of otherwise matching configuration; its initial weights were obtained by averaging the pretrained RGB filters across their input-channel dimension. The remaining feature-extraction layers were retained, while the final classifier linear layer was replaced by a layer with two outputs.

The model produced logits for the Good (class 0) and Artifact (class 1) categories. A softmax operation converted these logits to class probabilities, and the artifact probability was taken from the class-1 output. The classifier therefore estimated image-quality evidence from the localized brain crop rather than from a complete three-dimensional dynamic.

\subsection{Classifier Training and Model Selection}
\label{subsec:2d_classifier_training}

Classification targets were inherited from dynamic-level quality labels, and the subject-separated populations are described in Section~\ref{sec:data-preparation}. Separate classifier variants were trained with oracle and predicted crops.

Separate experimental variants used reference-mask-assisted (oracle) crops and automatically predicted crops. An early oracle-crop baseline was trained using class-weighted cross-entropy and the AdamW optimizer. The recorded configuration specified a maximum of 50 epochs, a learning rate of $10^{-4}$, weight decay of $10^{-4}$, early-stopping patience of eight epochs, and automatic mixed precision. The class weights were 1.0 for Good and 4.6174 for Artifact, reflecting the imbalance in the training labels. Model selection for this baseline was based on validation precision--recall performance.

Additional experiments were conducted with pretrained-named oracle and predicted-crop classifiers, and a predicted-crop checkpoint was used in the sequential demonstration. The available implementation verifies that the EfficientNet-B0 builder can initialize ImageNet weights and that saved classifier checkpoints contain training metadata, but does not independently establish the exact optimization configuration or checkpoint provenance for every later variant. The early oracle configuration was therefore not assumed to apply unchanged to all experiments.

\subsection{Slice-Level Artifact Probability Estimation}
\label{subsec:2d_slice_probability}

During classification inference, each available brain-centered crop was checked for the expected $224\times224$ spatial dimensions, converted to a single-channel tensor, and passed to the trained EfficientNet-B0 in evaluation mode. Nonfinite pixel values were sanitized and crop intensities were clipped to $[0,1]$. The inference path did not replicate grayscale images into RGB channels or apply three-channel ImageNet input normalization.

For logits $z_0$ and $z_1$ corresponding to Good and Artifact, respectively, the estimated probability of an artifact was calculated by
\begin{equation}
p_{\mathrm{artifact}}=
\frac{\exp(z_1)}{\exp(z_0)+\exp(z_1)}.
\end{equation}
The continuous artifact probability was retained for subsequent thresholding and dynamic-level analysis. Class probabilities were not computed when brain localization failed to generate a usable crop. Such missing classifier outputs were tracked separately from low predicted artifact probabilities.

For slice $s$ of dynamic $d$, denote the artifact-class score by $p_{d,s}$. A configurable operating threshold $\tau$ defines the binary decision
\begin{equation}
\hat y_{d,s}(\tau)=
\begin{cases}
1,&p_{d,s}\geq\tau,\\
0,&p_{d,s}<\tau.
\end{cases}
\end{equation}
These scores represent evidence for inherited dynamic-quality labels rather than independently calibrated probabilities of visible slice artifacts. Available slice probabilities, crop failures, and source identifiers were retained for the subsequent dynamic-level decision procedures.

\section{Dynamic-Level and Sequential Quality Assessment}
\label{sec:dynamic_sequential_assessment}

The 2D framework produced artifact-class probabilities for individual brain-centered MRI slices, whereas its source quality labels described complete diffusion dynamics. A further decision layer was therefore needed to combine slice outputs into acquisition-level assessments. Two distinct procedures were investigated: retrospective aggregation of slice-level scores (optionally incorporating crop failures) and a sequential demonstration in which the application could stop at an individual-slice trigger. The latter was an offline processing workflow and was not treated as a validated scanner-time early-detection result.

\subsection{Aggregation of Slice-Level Predictions}
\label{subsec:2d_dynamic_aggregation}

Let $\mathcal{S}^{\mathrm{valid}}_d$ denote the slice indices of dynamic $d$ for which segmentation, crop extraction, and classification succeeded. The available artifact-class scores were
\begin{equation}
\mathcal{P}_d=\{p_{d,s}:s\in\mathcal{S}^{\mathrm{valid}}_d\},
\qquad n_d=|\mathcal{S}^{\mathrm{valid}}_d|.
\end{equation}
For $n_d>0$, an aggregation operator $g$ produced a dynamic-level score $A_d=g(\mathcal{P}_d)$.

The recovered evaluation considered arithmetic mean, median, and maximum aggregation:
\begin{align}
A_d^{\mathrm{mean}}&=\frac{1}{n_d}\sum_{s\in\mathcal{S}^{\mathrm{valid}}_d}p_{d,s},\\
A_d^{\mathrm{median}}&=\operatorname{median}(\mathcal{P}_d),\\
A_d^{\mathrm{max}}&=\max_{s\in\mathcal{S}^{\mathrm{valid}}_d}p_{d,s}.
\end{align}
The mean summarized the overall slice evidence but could dilute isolated high scores. The median reduced the influence of extreme predictions, while the maximum preserved the strongest observed artifact score but could be sensitive to a single false positive.

Additional recorded configurations were identified as \texttt{top2}, \texttt{top3}, and \texttt{top5}. These labels referred to variants selecting high-scoring slices, but the recovered outputs did not conclusively establish the calculation performed for every archived configuration. An explicit top-$k$ equation was therefore not attributed to all such results without the corresponding executed script.

Slices lacking a usable crop contributed no classifier probability; their missing scores were not replaced with zero. When no valid scores were available, aggregation was undefined and required a separate failure policy. Each dynamic, rather than each constituent slice, was treated as one evaluation observation for acquisition-level quality assessment.

\subsection{Artifact Classification Thresholds}
\label{subsec:2d_dynamic_thresholds}

A decision threshold converted the aggregated Artifact-class score to a binary quality prediction:
\begin{equation}
D_d^{\mathrm{cls}}(\tau)=
\begin{cases}
1,& A_d\geq\tau,\\
0,& A_d<\tau.
\end{cases}
\end{equation}
Class 1 indicated Artifact and class 0 indicated Good, under the selected aggregation operator and threshold. Lower thresholds could increase sensitivity at the cost of false alerts; higher thresholds could reduce false positives while missing artifact-positive acquisitions. Thresholds were therefore analyzed jointly with the aggregation configuration and evaluation cohort.

Offline retrospective analyses explored combinations of aggregation operators and candidate thresholds, including failure-aware settings. Their purpose was to characterize sensitivity, specificity, precision, and other dynamic-level metrics across operating points. An exploratory threshold-search result did not automatically constitute an operating rule frozen before independent test evaluation, particularly when complete selection provenance was not available.

The sequential demonstration used a separate user-adjustable \emph{slice-level} threshold initialized at 0.92, with an allowed range of 0.50--0.99. That condition compared one current slice score with its threshold; it was not equivalent to thresholding a full-dynamic mean, median, maximum, or top-$k$ aggregation. The experimental settings were therefore kept distinct.

\subsection{Segmentation-Failure-Aware Decision Rules}
\label{subsec:2d_failure_aware_rules}

The classifier required a valid brain crop. An empty predicted mask or unsuccessful crop prevented ordinary classification and was recorded as a localization or crop-generation failure. Conversely, a nonempty predicted mask could yield a crop even when its anatomical overlap was poor; this was not the same event as a missing crop.

Define the observed failure indicator for selected slice $s$ of dynamic $d$ as
\begin{equation}
F_{d,s}=
\begin{cases}
1,&\text{no usable brain crop is generated},\\
0,&\text{a usable crop is generated}.
\end{cases}
\end{equation}
For selected slice set $\mathcal{S}_d$, the failure count was
\begin{equation}
N_d^{\mathrm{fail}}=\sum_{s\in\mathcal{S}_d}F_{d,s}.
\end{equation}
A recovered failure-aware dynamic assessment combined an aggregated classifier score with a minimum failure-count condition. For a valid aggregated score, its decision may be expressed as
\begin{equation}
D_d^{\mathrm{QC}}=
\mathbb{1}\!\left[A_d\geq\tau\ \lor\ N_d^{\mathrm{fail}}\geq\kappa\right],
\end{equation}
where $\kappa$ denotes the configured failure threshold. The recovered evaluation included a setting with $\kappa=1$. This expression describes the documented failure-aware policy family; the exact semantics of all archived failure-threshold labels could not be established from the output tables alone. Cases without any valid classification scores must be handled directly by a specified failure policy, rather than by evaluating an undefined $A_d$.

A failure-triggered flag was a quality-control decision, not an independently confirmed artifact. An unusable mask could reflect image degradation, anatomy, slice selection, or segmentation-model limitations. Failure-aware evaluation therefore retained failures in the appropriate dynamic-level denominators and distinguished classifier-only performance among usable crops from behavior of the entire processing workflow.

\subsection{Sequential Processing and Early-Stopping Logic}
\label{subsec:2d_sequential_stopping}

A sequential demonstration accepted reconstructed 3D or 4D NIfTI input, selected one dynamic where required, and processed an ordered range of image slices. The documented default used the third image axis and indices 8--25, inclusive, without an additional in-plane transformation. For each slice, the application performed image normalization, MobileNetV3-Large U-Net inference, restoration of the predicted binary mask to source dimensions, brain-centered crop extraction, and EfficientNet-B0 artifact-probability estimation when possible. Slice outputs, failure states, and component timings were recorded as processing progressed. Model checkpoints were cached rather than reloaded for each slice.

\paragraph{Artifact-trigger stopping.}
For current slice $s$, a usable crop produced probability $p_{d,s}$. The application compared it with the slice-level threshold $\tau_s$, generating an artifact trigger when
\begin{equation}
T_{d,s}^{\mathrm{artifact}}=\mathbb{1}[p_{d,s}\geq\tau_s]=1.
\end{equation}
If this condition was reached, the application flagged the dynamic and stopped examining subsequent selected slices.

\paragraph{Optional failure-trigger stopping.}
When the stop-on-failure option was enabled, an empty predicted mask or unusable crop could also terminate the sequence with a distinct failure designation. With $\eta\in\{0,1\}$ indicating whether that option was active, the combined trigger was
\begin{equation}
T_{d,s}=T_{d,s}^{\mathrm{artifact}}\lor(\eta F_{d,s}),
\end{equation}
with the classifier trigger treated as false when no classifier prediction existed. This logical rule should not be confused with the retrospective failure-count threshold $\kappa$.

For an ordered set of selected slices $s_1,\ldots,s_K$, the first-trigger position was defined, if a trigger occurred, as
\begin{equation}
J_d=\min\{j:T_{d,s_j}=1\}.
\end{equation}
The decision position counted examined slices rather than representing the original absolute slice index. In the absence of a trigger, processing reached the final selected slice; this outcome meant only that no enabled stopping condition occurred within that range, not that the entire acquisition was demonstrably artifact-free.

The demonstration implemented logical early stopping over slices from a completed input image. It did not establish prospective streaming from the scanner or automatic reacquisition. The recovered evidence did not provide a verified per-dynamic first-trigger evaluation or measured time saved through early stopping. Arithmetic estimates based on unprocessed slice counts were therefore not treated as actual observed latency reductions.

\subsection{Dynamic-Level Quality Decisions}
\label{subsec:2d_final_dynamic_decisions}

The developed decision layer distinguished retrospective dynamic classification from sequential first-trigger quality-control outcomes.

\paragraph{Retrospective dynamic decisions.}
Available slice probabilities were grouped by original subject and dynamic identifiers. An aggregation operator generated a score and a selected threshold produced the classifier-based binary decision. Failure-aware configurations additionally considered the number of unusable crops. Each evaluated dynamic contributed one reference quality label and one resulting decision for a specified rule, allowing assessment using accuracy, Artifact sensitivity, specificity, precision, F1-score, and corresponding confusion-matrix counts.

\paragraph{Sequential demonstration outcomes.}
The ordered application distinguished (i) an artifact-related trigger based on a classifier score, (ii) a localization-failure trigger when enabled, and (iii) completion of the selected range without a trigger. This separation was necessary because a failed segmentation was not a classifier-confirmed artifact. The sequential decision did not require all selected slice probabilities, unlike retrospective full-dynamic aggregation.

\paragraph{Interpretation and evaluation boundary.}
Retrospective dynamic results establish agreement of recorded decisions with acquisition-level annotations for their specific test populations and operating conditions. Demonstrating useful early detection additionally requires validated first-decision positions, the decision information available before each trigger, reliable ground truth, and measured processing time to that decision. Such a complete sequential evaluation was not established in the recovered results. The broader experimental configurations and metric definitions are described in Section~\ref{sec:experimental-evaluation}, with numerical findings presented in Chapter 5.

\section{Experimental Design and Evaluation Protocol}
\label{sec:experimental-evaluation}

The experimental evaluation examined segmentation quality, artifact classification, reliability of brain-centered cropping, dynamic-level decisions, and computational feasibility. Because the research included preliminary 3D experiments, a developed 3D segmentation-to-classification pipeline, and a subsequent 2D framework, evaluation populations and prediction units differed across experiments. These differences were maintained when interpreting the outcomes. Segmentation predictions were compared with spatial fetal-brain masks; 2D slice-classification predictions were evaluated against quality labels inherited from their parent dynamics; and dynamic-level quality decisions were evaluated against the original acquisition-level annotations.

\subsection{Experimental Configurations}
\label{subsec:experimental-configurations}

\paragraph{Preliminary volumetric experiments.}
The initial experiments combined a custom 3D U-Net with brain cropping, handcrafted statistical features, and Random Forest classification. The exploratory Random Forest experiment contained four feature records and used an incorrect subject-to-quality-label association. Its numerical performance was therefore not used as a valid independent classification benchmark or to rank models. The custom U-Net's dynamic-entry partitions were not subject independent.

\paragraph{Developed volumetric pipeline.}
The developed 3D pipeline paired nnU-Net v2 brain segmentation with 3D ResNet artifact classification. The documented nnU-Net preparation comprised 600 training and 170 test image--mask entries, with some subject identities occurring in both partitions. The recovered classifier data comprised 600 eligible dynamics with recorded predictions, but the training exposure of these observations and the precise checkpoint-to-export provenance were not conclusively established. Consequently, results reconstructed from this export were not assumed to represent an independently held-out test evaluation.

\paragraph{Developed slice-wise pipeline.}
The 2D approach combined MobileNetV3-Large U-Net segmentation, segmentation-dependent ROI cropping, and EfficientNet-B0 classification. Its recorded cohort contained 24 subject groups assigned to training, validation, and testing as 16, four, and four groups, respectively. The documented classification metadata supported disjoint subject identities across partitions. The original split-generation procedure and complete identifier mapping were not independently verified in every recovered record. Brain segmentation and classification used different prepared slice populations, with the segmentation data retaining background-only slices and classification restricted to selected brain-containing slices. Counts and selection criteria are defined in Section~\ref{sec:data-preparation}.

\paragraph{Quality-decision experiments.}
Retrospective experiments investigated aggregation of slice probabilities, operating thresholds, and segmentation-failure-aware dynamic decisions. A separate sequential demonstration iterated over selected slices and applied individual-slice and optional failure-trigger stopping conditions. The retrospective results and demonstration were treated as different protocols; a retrospective dynamic-level score did not establish a prospectively achieved early decision.

Across experiments, observed results were separated from unresolved provenance and from hypothetical follow-on tests. Numerical outcomes are presented in Chapter 5.

\subsection{Segmentation Evaluation Protocol}
\label{subsec:segmentation-evaluation}

The developed 2D segmentation network was evaluated against reference binary fetal-brain masks in the subject-separated validation and test partitions. Segmentation logits were converted to foreground probabilities using sigmoid activation, followed by thresholding at 0.5. Brain pixels were positive and background pixels negative.

For pixel-level true positives ($TP$), false positives ($FP$), false negatives ($FN$), and true negatives ($TN$), the evaluated metrics were
\begin{align}
\mathrm{Dice} &= \frac{2TP}{2TP+FP+FN}, &
\mathrm{IoU} &= \frac{TP}{TP+FP+FN},\\
\mathrm{Sensitivity} &= \frac{TP}{TP+FN}, &
\mathrm{Specificity} &= \frac{TN}{TN+FP}.
\end{align}
A numerical stabilizer of $10^{-7}$ was used in implemented calculations where needed to avoid division by zero. The documented evaluation computed measurements per slice before averaging them. Therefore, aggregate mean Dice and mean IoU need not obey the exact relationship that holds for one individual binary mask or a single pooled confusion matrix.

The evaluation population included slices without reference brain foreground. Under an empty-mask smoothing convention, an empty reference paired with an empty prediction may receive perfect overlap even though no brain has been localized. Brain-present and background-only slices were therefore distinguished when interpreting overlap measures. Predicted masks were also examined for empty foreground and poor anatomical agreement, since successful crop generation and accurate brain localization are distinct outcomes.

Segmentation performance was examined by the parent dynamic's inherited Good or Artifact label for selected slices. These subgroup comparisons measured an association between acquisition quality annotations and segmentation behavior; they did not assume that an artifact was visually present in every individual positively labelled slice. For the final nnU-Net stage, a complete verifiable held-out set of predictions, reference masks, and aggregate segmentation scores was not recovered, and independently validated segmentation performance was not inferred from downstream classifier availability.

\subsection{Artifact Classification Evaluation Protocol}
\label{subsec:artifact-classification-evaluation}

For classification, Good was class 0 and Artifact class 1. The 3D classifier generated one prediction for an eligible dynamic, whereas the 2D classifier generated a prediction for each usable brain-centered slice. In the latter case, the reference label was inherited from the parent dynamic, constituting weak supervision rather than expert annotation of visible slice-level corruption.

At a selected threshold, true positives were correctly flagged Artifact-labelled observations; false positives were Good-labelled observations incorrectly flagged as Artifact; false negatives were Artifact-labelled observations classified as Good; and true negatives were correctly classified Good-labelled observations. The classification measures were
\begin{align}
\mathrm{Accuracy} &= \frac{TP+TN}{TP+TN+FP+FN},\\
\mathrm{Precision} &= \frac{TP}{TP+FP}, &
\mathrm{Recall} &= \frac{TP}{TP+FN},\\
\mathrm{Specificity} &= \frac{TN}{TN+FP}, &
F_1 &= \frac{2TP}{2TP+FP+FN}.
\end{align}
Undefined denominator cases require an explicitly documented numerical convention, rather than implicit interpretation.

Receiver operating characteristic area under the curve (ROC-AUC) and precision--recall analysis were used to examine discrimination over thresholds. Where average precision was calculated, it was distinguished from trapezoidal precision--recall AUC, as these are not identical numerical summaries. Sensitivity and precision were emphasized alongside overall accuracy because Good labels predominated.

Classifier-only metrics were calculated on observations with available model predictions. Missing crops did not receive an artificial low Artifact probability and were evaluated separately through crop availability and failure-aware analysis. The analysis unit, class counts, decision threshold, and population denominator were retained for each reported result.

\subsection{Oracle Versus Predicted Crop Comparison}
\label{subsec:cropping-comparison}

Two brain-region localization conditions were examined. Oracle cropping used reference brain masks; predicted cropping used MobileNetV3-Large U-Net outputs. Both conditions used the same 15\% bounding-box margins, square-crop construction, boundary padding, intensity preparation, and $224\times224$ resizing described in Section~\ref{subsubsec:roi-cropping}. Reference-assisted localization was not assumed to provide perfect anatomical ground truth.

The documented test population contained 1,529 eligible reference-mask crops. Predicted cropping provided usable inputs for 1,499 of these slices; 30 had no valid predicted crop. The matched classification comparison therefore used the 1,499 observations available under both conditions, preserving their common subject, dynamic, slice, and label identities. Crop failures were retained in the separate ROI-reliability analysis and were not treated as ordinary Good classifications.

The recorded oracle and predicted classification variants were trained separately. Their matched evaluation consequently compared complete training-and-inference configurations, rather than isolating the causal effect of the mask source while holding the classifier weights fixed. Applying one frozen classifier to paired oracle and predicted crops would offer a more controlled measure of crop-input effects, but that specific controlled comparison was not established from the recovered results.

Classification metrics were examined alongside predicted-crop availability and segmentation overlap. A nonempty predicted mask could yield a crop despite poor anatomical agreement; thus, technical crop success and localization accuracy were examined separately.

\subsection{Dynamic-Level Evaluation}
\label{subsec:dynamic-evaluation}

Dynamic-level analysis combined slice-level classifier outputs using subject and dynamic identifiers, yielding one quality score or decision for each evaluated dynamic. The original dynamic-level Good or Artifact annotation served as the reference target. Mean, median, maximum, and recorded top-$k$ aggregation variants were evaluated as described in Section 4.7. The archived top-$k$ method labels do not independently establish the exact mathematical implementation for every historical evaluation.

Candidate probability thresholds were investigated using confusion matrices and class-sensitive metrics. Failure-aware evaluation additionally incorporated unavailable segmentation-derived crops, distinguishing classifier outputs from missing classifier inputs. The recovered implementation included a decision rule that flagged a dynamic when either the aggregated score reached its threshold or the number of failed crops reached a selected failure-count threshold. These operating-point searches were not assumed to have been frozen independently of all later test inspection unless execution provenance established that separation.

The archived validation-based aggregation search contained 86 dynamics, with 80 Good and six Artifact labels. A separate test evaluation contained 88 dynamics, with 70 Good and 18 Artifact labels. Their populations, thresholds, and denominators were kept separate from other experiments.

Dynamic-level metrics included the confusion matrix, accuracy, Artifact sensitivity, specificity, precision, and F1-score, with ROC-AUC and precision--recall summaries where continuous dynamic scores were available. Each dynamic contributed one reference quality label and one decision per evaluated configuration. Slices from the same dynamic were not treated as independent dynamic-level observations.

The sequential demonstration used a different evaluation mechanism: ordered slice inference and a first-trigger stopping rule. The recovered records did not establish a verified quantitative evaluation of first-trigger positions, early-detection outcomes using only observations available before each decision, or actually measured saved time. Accordingly, retrospective aggregation results were not described as prospectively validated early stopping.

\subsection{Matched 3D Versus 2D Comparison}
\label{subsec:matched-3d-2d-evaluation}

The two developed pipelines were compared on a matched set of diffusion dynamics for which recorded quality decisions from both approaches were available. The population contained 71 dynamics from four subjects, comprising 59 Good and 12 Artifact-labelled acquisitions. Both methods were evaluated against the same dynamic-level ground-truth annotations on the same 71 identities.

The 3D approach supplied volumetric ResNet-based quality predictions. The 2D approach supplied recorded dynamic-level decisions derived from slice-level evidence and a segmentation-failure indicator. True positives, false positives, false negatives, and true negatives, together with accuracy, Artifact sensitivity, specificity, precision, and F1-score, were compared on the matched population.

Although matched identities remove differences in the evaluated observations, they do not establish equivalent training-data independence. The 3D classifier's training exposure to the matched population could not be conclusively excluded, and original execution provenance and operating-rule selection details were not verified in all respects. The comparison was consequently interpreted as an exploratory matched analysis of recovered predictions, not definitive evidence of generalization to independent fetuses.

The matched sample contained only four subjects, so dynamics within a subject were not independent biological samples. Formal paired significance testing and subject-clustered uncertainty intervals had not been established in the documented analysis. These limitations were retained when interpreting apparent model differences.

\subsection{Computational Performance Measurement}
\label{subsec:computational-evaluation}

Computational evaluation investigated the timing of image preparation, segmentation inference, brain-crop extraction, artifact classification, and quality-decision processing. Distinct measurement scopes were maintained when considering 2D inference benchmarks, prepared-volume classification, and earlier offline 3D workflow simulation.

\paragraph{Two-dimensional measurements.}
The 2D timing protocol examined segmentation preprocessing, MobileNetV3-Large U-Net inference, ROI preparation, EfficientNet-B0 inference, and aggregation where recorded. A recovered batch-size-one segmentation benchmark used 50 warm-up iterations and 200 measured iterations, with automatic mixed precision. Stage-specific measurements were summarized using means, medians, upper percentiles, and maxima where available. The population of candidate slices was distinguished from the subset with successful predicted crops and classifier inputs.

\paragraph{Volumetric measurements.}
The recovered 3D classifier benchmark used already prepared volumetric crops and did not represent the full nnU-Net-to-classification path. Historical folder-monitoring simulation reports covered a broader processing sequence but were not demonstrated to share the same hardware, checkpoint versions, or timing boundaries as later benchmarks. Therefore, 3D classifier-only times and complete 2D slice-processing times were not directly ranked as equivalent end-to-end measures.

\paragraph{Sequential timing interpretation.}
The sequential demonstration recorded processing steps and stop-trigger conditions. Actual time to a decision would require a measured first-trigger time and an appropriate no-stopping comparator. Multiplying the count of unprocessed slices by mean per-slice inference time yields only a projection, not an independently measured acquisition-time saving. Complete scanner-to-decision latency was not established.

Disk input/output, checkpoint loading, GPU synchronization, host--device transfer, visualization delays, scanner acquisition, and control-system communication were not assumed to be included unless measurement records specified them. Reported timing conclusions were limited to the instrumented stages.

\subsection{Software and Hardware Environment}
\label{subsec:software-hardware-environment}

The experimental workflows were implemented in Python. NiBabel was used to load NIfTI MRI images and masks, NumPy for array manipulation and intensity processing, and pandas for annotation tables and prediction records. PyTorch supported custom segmentation and classification training and inference, while torchvision provided model components for MobileNetV3-Large and EfficientNet-B0. The volumetric segmentation stage additionally used nnU-Net v2. The exploratory Random Forest model, stratified splitting, and classification evaluation used scikit-learn.

Matplotlib supported diagnostic visualization and training/evaluation figures; ITK-SNAP was used where applicable for inspection of volumetric data and masks. Streamlit provided interactive processing demonstrations, including the slice-wise sequential application.

CUDA acceleration and automatic mixed precision were used where compatible resources and scripts permitted. The recovered 2D latency benchmark identified an NVIDIA RTX PRO 6000 Blackwell Server Edition GPU partition configured as MIG 1g.24gb. This hardware was not assumed to represent all historical training and inference runs. Likewise, fixed seeds documented for selected procedures did not establish complete reproducibility across unspecified execution environments. Full run-specific package versions, hardware inventories, and checkpoint provenance could not be reconstructed for every experiment.

Overall, the evaluation protocol distinguished anatomical segmentation from quality classification, slice-level weak targets from dynamic-level reference annotations, classifier-only performance from crop and processing failures, and offline computation from acquisition-time operation. These distinctions govern the interpretation of the experimental findings reported in Chapter 5.

